\documentclass[12pt]{article}
\usepackage[utf8]{inputenc}
\usepackage[T1]{fontenc}
\usepackage[margin=1in]{geometry}
\usepackage{setspace}
\usepackage{xcolor}
\usepackage[colorlinks=true, urlcolor=blue!60!black, linkcolor=blue!60!black]{hyperref}
\setlength{\parindent}{0em}
\setlength{\parskip}{0.6em}
\renewcommand{\baselinestretch}{1.0}
\begin{document}
\begin{center}
{\LARGE \textit{Seminole Tribe of Florida v.\ Florida}, 517 U.S.\ 44 (1996)\par}
\vspace{0.5em}
{\large Case Brief --- Two-Page Edition\par}
\end{center}

\section{Memory Jogger}
Congress cannot use Article I to abrogate state sovereign immunity, and \textit{Ex parte Young} could not bypass IGRA's remedial scheme (\href{https://www.law.cornell.edu/supct/html/94-12.ZO.html}{Cornell LII}).

\section{Detailed Case Facts}
The Indian Gaming Regulatory Act (IGRA), enacted under the Indian Commerce Clause, required States to negotiate gaming compacts with tribes in good faith and let tribes sue States in federal court, followed by a mediation and Secretarial-procedures ladder. The Seminole Tribe sued Florida and Governor Chiles for refusing to negotiate.

\section{Procedural History}
The district court upheld abrogation; the Eleventh Circuit reversed and ordered dismissal; the Supreme Court affirmed.

\section{Judicial Votes}
Affirmed, 5--4. Rehnquist wrote for O'Connor, Scalia, Kennedy, and Thomas; Stevens dissented, as did Souter, joined by Ginsburg and Breyer.

\section{Holding}
IGRA clearly sought to abrogate immunity, but \textit{Union Gas} is overruled: Article I powers cannot abrogate the immunity of unconsenting States, though Section~5 of the Fourteenth Amendment can. A detailed statutory remedial scheme displaces \textit{Ex parte Young} suits against state officers.

\section{Analysis of Opinions}
\textbf{Majority:} Under \textit{Hans}, immunity is a constitutional ``presupposition'' that protects States from ``the indignity'' of private suits; \textit{Fitzpatrick} rested on the later Fourteenth Amendment, and \textit{Union Gas} drew no principled line. \textbf{Stevens:} the ruling is an ``affront to a coequal branch'' leaving federal rights without remedy. \textbf{Souter:} the Eleventh Amendment repealed only diversity jurisdiction; elevating common-law immunity above Article I repeats \textit{Lochner}; \textit{Young} is ``indispensable.''

\section{Examples: Future Applications}
\textbf{Same-Side:} patent or copyright damages suits against States under Article I are barred. \textbf{Opposite-Side:} family-leave damages suits valid under Section~5 (\textit{Hibbs}) proceed. \textbf{Fence-Sitter:} abrogation under the Bankruptcy Clause or war powers, where States may have surrendered immunity in ``the plan of the convention.''

\section{Critique}
Critics note the irony of textualists endorsing an atextual immunity, though Baude and Sachs defend it as a Founding-era backdrop. Vicki Jackson and Erwin Chemerinsky warn of a federalism pincer with \textit{Boerne} and of weakened \textit{Ex parte Young} enforcement.

\section{Key Quotations}
\begin{quote}
``Even when the Constitution vests in Congress complete lawmaking authority over a particular area, the Eleventh Amendment prevents congressional authorization of suits by private parties against unconsenting States. The Eleventh Amendment restricts the judicial power under Article III, and Article I cannot be used to circumvent the constitutional limitations placed upon federal jurisdiction.''
\end{quote}
\noindent --- Chief Justice Rehnquist, Opinion of the Court.

\end{document}
